\documentclass[11pt]{article}
\usepackage[a4paper,margin=1in]{geometry}
\usepackage{fontspec}
\usepackage[boldfont=false]{xeCJK}
\setCJKmainfont{FandolSong-Regular.otf}
\usepackage{amsmath}
\usepackage{amssymb}
\usepackage{array}
\usepackage{booktabs}
\usepackage{graphicx}
\usepackage{adjustbox}
\usepackage{enumitem}
\setlist[itemize]{topsep=2pt,partopsep=0pt,parsep=0pt,itemsep=2pt,leftmargin=1.4em}
\setlist[enumerate]{topsep=2pt,partopsep=0pt,parsep=0pt,itemsep=2pt,leftmargin=1.6em}
\usepackage{parskip}
\usepackage{fvextra}
\fvset{breaklines=true,breakanywhere=true,fontsize=\small}
\setlength{\parindent}{0pt}

\begin{document}

\begin{center}
  {\LARGE\bfseries Genetic technology}\\[0.35em]
  {\large A Level Biology}
\end{center}
\vspace{0.6em}

\section*{Principles of genetic technology}

\textbf{Recombinant} 重组 DNA is DNA that has been made by joining together DNA from two different sources. \textbf{Genetic engineering} 基因工程 is the deliberate changing of an organism's genetic material — often by transferring a \textbf{gene} 基因 into an organism so that the gene is \textbf{expressed} (switched on to make its \textbf{protein} 蛋白质).

The gene to be transferred can be obtained in three ways:

\begin{itemize}
\item cut out of the DNA of a \textbf{donor} 供体 organism,

\item made from the donor's mRNA, using the enzyme \textbf{reverse transcriptase} 逆转录酶,

\item built chemically from \textbf{nucleotides} 核苷酸.

\end{itemize}

\subsection*{The tools}

\begin{center}
\adjustbox{max width=\linewidth}{%
\begin{tabular}{ll}
\toprule
\textbf{Tool} & \textbf{Role} \\
\midrule
\textbf{restriction endonuclease} 限制性内切酶 & cuts DNA at a specific base sequence, leaving "sticky ends" \\
DNA \textbf{ligase} 连接酶 & joins pieces of DNA together \\
\textbf{plasmid} 质粒 & a small ring of bacterial DNA used as a \textbf{cloning vector} 载体 to carry the gene into a cell \\
DNA \textbf{polymerase} 聚合酶 & copies DNA \\
reverse transcriptase & makes DNA from an mRNA template \\
\bottomrule
\end{tabular}}
\end{center}

A \textbf{promoter} 启动子 often has to be transferred along with the gene. The promoter is the "switch" that lets the gene be transcribed in its new organism, so without it the gene would stay silent.

To check the gene has gone in and is working, scientists add a \textbf{marker gene} 标记基因 next to it — for example one that codes for a \textbf{fluorescent} 荧光 (glowing) product. If the cells glow, the transfer worked.

\subsection*{Gene editing}

\textbf{Gene editing} 基因编辑 is a precise form of genetic engineering. It inserts, deletes or replaces DNA at an exact site in the \textbf{genome} 基因组.

\subsection*{Copying and sorting DNA}

\begin{itemize}
\item the \textbf{polymerase chain reaction} 聚合酶链式反应 (PCR) is used to \textbf{clone} 克隆 and \textbf{amplify} 扩增 DNA — to make millions of copies. It repeats cycles of heating and cooling, using a heat-stable enzyme called Taq polymerase.

\item \textbf{gel electrophoresis} 凝胶电泳 separates DNA \textbf{fragments} 片段 by length. The fragments move through a gel in an electric field, and shorter fragments move further, so the lengths spread out into bands.

\item \textbf{microarrays} 微阵列 are used to study whole \textbf{genomes} and to detect which genes are switched on, by picking up their mRNA.

\item \textbf{databases} 数据库 store the nucleotide sequences of genes and the \textbf{amino acid} 氨基酸 sequences of proteins, so scientists anywhere can compare them.

\end{itemize}

\section*{Genetic technology in medicine}

\subsection*{Recombinant human proteins}

A human gene can be put into bacteria or other cells so that they make a \textbf{recombinant} human \textbf{protein} — an exact copy of the human one. This is safer and never in short supply, and it avoids using proteins taken from animals or donors. Examples are \textbf{insulin} 胰岛素 (for \textbf{diabetes} 糖尿病), factor VIII (for haemophilia) and adenosine deaminase (for a faulty immune system).

\subsection*{Genetic screening}

\textbf{Genetic screening} 基因筛查 tests a person's DNA for disease alleles before symptoms appear. Examples are the BRCA1 and BRCA2 alleles (which raise the risk of \textbf{breast cancer} 乳腺癌), \textbf{Huntington's disease} 亨廷顿病, and \textbf{cystic fibrosis} 囊性纤维化. Knowing the result helps people make informed choices about treatment and family.

\subsection*{Gene therapy}

\textbf{Gene therapy} 基因治疗 treats a genetic disease by putting a working copy of a gene into the patient's cells. It has been used for SCID (a disease in which the immune system fails) and for some inherited eye diseases.

\subsection*{Social and ethical questions}

Genetic screening and gene therapy raise concerns: who should see your genetic results, whether insurers or employers could misuse them, whether changes are safe and permanent, and who decides. These \textbf{ethical} 伦理 and social questions must be weighed against the benefits.

\section*{Genetically modified organisms in agriculture}

\textbf{Genetic engineering} can help feed a growing world by improving farmed animals and crops. Examples of \textbf{genetically modified organisms} 转基因生物 (GMOs) are:

\begin{itemize}
\item \textbf{GM salmon} 鲑鱼 that grow to size faster.

\item \textbf{soybean} 大豆 made resistant to a \textbf{herbicide} 除草剂, so weeds can be sprayed without harming the crop.

\item \textbf{cotton} 棉花 made resistant to insect \textbf{pests} 害虫, because it makes a protein that kills the insects.

\end{itemize}

GMOs also raise ethical and social questions: whether they are safe to eat, what effect they have on the environment and wild species, whether the engineered genes might spread, and whether a few large companies should control the food supply.


\end{document}
